\uuid{Bhsd}
\exo7id{3423}
\titre{exo7 3423}
\auteur{quercia}
\organisation{exo7}
\datecreate{2010-03-10}
\isIndication{false}
\isCorrection{false}
\chapitre{Déterminant, système linéaire}
\sousChapitre{Système linéaire, rang}
\cpge{mpsi-mp2i.matrices-applications-lineaires.b, pcsi.matrices-applications-lineaires.b, ptsi.matrices-applications-lineaires.b, mpsi-mp2i.matrices-applications-lineaires.c, pcsi.matrices-applications-lineaires.d, ptsi.matrices-applications-lineaires.d, mpsi-mp2i.determinants.d, pcsi.calcul-matriciel-systemes-lineaires.d, ptsi.calcul-matriciel-systemes-lineaires.d, bcpst1.systemes-lineaires.a, bcpst1.matrices, pc.complements-espaces-vectoriels.e, psi.complements-espaces-vectoriels.e, bcpst2.applications-lineaires-matrices.a, bcpst2.applications-lineaires-matrices.c}
\module{Algèbre}
\niveau{L2}
\difficulte{}

\contenu{
\question{
Soient $\alpha_1, \dots, \alpha_n$, $n$ scalaires distincts et
$M$ la matrice $\big(\alpha_i^{j-1}\big)$ (matrice de Vandermonde).


Montrer que $M$ est inversible en interprétant le système $MX = 0$ dans
$ K_{n-1}[x]$.
}
}
